\section{Conclusions, limitations and plans}
\label{sec:conclusions}

\NHFOAM\ follows NHWAVE's pressure arrangement on OpenFOAM's unstructured,
moving $\sigma$-meshes. The tests of \citet{ma2012} were run with both models
on the same set-ups, from the simplest upwards, and they changed the code:
\begin{itemize}[nosep]
  \item \textbf{Time accuracy.} With one outer corrector the convection term
        was linearised with the flux of the previous time level, which made the
        scheme first order in time for nonlinear waves. The convecting flux is
        now extrapolated from the last two levels. A solitary wave of
        $H/h=0.5$ lost 6.4\% of its height in 150 m before and 1.0\% after the
        fix (NHWAVE 2.3\%), at no extra cost.
  \item \textbf{Convection.} The TVD scheme damps steep short waves even when
        they do not break; filtered central differencing
        (\code{filteredLinear2V 0.2 0}) is now recommended for non-breaking
        waves. It reduced the standing-wave error to 0.16--0.25\% for
        $kD\ge 2\pi$ (NHWAVE 0.6--4.8\%) and stayed stable in 40--65 s runs over
        the submerged bar. Breaking still needs the TVD scheme.
  \item \textbf{Consistency with NHWAVE.} After the fix, the phase of the waves
        over and behind the submerged bar agrees with NHWAVE's. The closer
        match of the measured phase before the fix came from the first-order
        lag, which slowed the waves by about 0.5\%.
  \item \textbf{Resolution.} On the elliptic shoal, with twice the paper's
        grid spacing and three layers, \NHFOAM's errors at the first sections
        (0.06 and 0.13) are close to NHWAVE's on the paper's grid (0.07 and
        0.11), while NHWAVE on the coarse grid loses 30\% of the incident wave
        height before the shoal.
  \item \textbf{Robustness.} A domain in which waves break anywhere needs the
        TVD scheme and the breaking model, even outside the region of
        interest: the shoal run with the filtered scheme failed when the waves
        broke in a shallow corner, and ran on with \code{limitedLinearV} and
        the hydrostatic front approximation.
  \item \textbf{Wetting and drying.} NHWAVE's mask and thin film carry over,
        but the implicit free surface needs a limiter on the outflows (which
        keeps the volume conserved, unlike NHWAVE's clip), and the collocated
        mesh needs the upwind (TVD) depth at the side faces and the upwind
        predicted flux in supercritical flow. With the centred depth of the
        mesh faces the fronts were too slow ($R/d = 0.41$ for the solitary
        wave); with both changes the run-up is 0.531 (NHWAVE 0.512, lab
        0.49--0.55), still water stays at rest on a beach, and the
        parabolic basin is reproduced to 0.9 mm rms. A dam break on a dry
        bed showed that both models treat the film as a wet bed and follow
        Stoker's solution for it, not Ritter's. Parallel runs give the same
        run-up with a processor boundary in the swash zone (0.540 against
        0.531).
  \item \textbf{Breaking.} The HFA alone flags only the fronts of the
        bores. A persistence threshold that falls with the age of the
        breaking event and the breaking eddy viscosity of
        \citet{kennedy2000}, with horizontal mixing, nearly halve the
        surf-zone errors of the wave height on both spilling-breaker beaches (mean
        error $+6.1 \to +3.5$ cm on the beach of \citet{ting1994}, NHWAVE
        $+4.8$; rms $2.98 \to 1.54$ cm on that of \citet{hansen1979}), raise
        the set-up, and keep the measured shape of the undertow, which is
        30\% too strong (NHWAVE 60\%). Vertical mixing through the layer
        faces flattens the undertow and dissipates less. In 3D, columns
        must join a breaker at the threshold of a new breaker, or the
        breaking region creeps along oblique crests into deep water.
  \item \textbf{Longshore current.} Test 4.7 of \citet{ma2012}, in 3D
        with periodic lateral boundaries, wetting and drying, log-law
        friction and the Smagorinsky model, runs with both models on the
        same coarse grid. Both put the current at about the measured
        position (\NHFOAM: 0.265 m/s at 1.08 m from the shoreline, NHWAVE:
        0.250 m/s at 1.26 m, lab 0.39 m/s at 1.21 m) and give two thirds of
        its strength. With the HFA alone the current is 0.19 m/s and too
        close to the shoreline, as \citet{ma2012} found without turbulent
        diffusion.
\end{itemize}

\paragraph{Limitations.}
\begin{itemize}[nosep]
  \item \textbf{Wetting and drying} has been tested in 2DV only, with uniform
        horizontal cells near the shoreline. The film acts as a wet bed, so
        thin swash flows run slower than over a truly dry bed (in NHWAVE as
        well), and puddles of a few millimetres above the film stay on the
        beach in the backwash. A narrow cell at the shoreline raises the
        errors. It needs the interface pressure layout.
  \item \textbf{Breaking} has been tested on two spilling-breaker beaches
        and the longshore current. With the eddy viscosity the decay inside
        the surf zone is close to the measured one, but just after the
        breaking point the models lose too little height (the heights are
        3.5--4 cm too large at $x = 8$ m on the beach of \citet{ting1994}),
        the breaking point is 0.2--0.6 m from the measured one, and the
        parameters
        ($\alpha$, $\beta$, $\beta_F$, $T_s$, $\delta_b$) are those of
        \citet{smit2013} and \citet{kennedy2000}, not tuned. The kinematic
        onset needs a better estimate of the crest speed than $\sqrt{gD}$.
  \item \textbf{Wave generation} uses relaxation zones with linear or
        stream-function targets. There is no boundary wavemaker and no
        irregular input.
  \item \textbf{Long waves in closed flumes.} The zones absorb long waves only
        partly and the wave-making zone does not conserve volume, so long waves
        of a few millimetres build up and shift the phase over shallow bars in
        long runs.
  \item \textbf{Restarts} do not save the time history of $\eta$ and
        $\bm{u}_f$; the first step after a restart is Euler and first order.
  \item \textbf{Landslide trailing waves} are 1.5--3 times the measured ones
        (Sections~\ref{sec:seabed} and \ref{sec:ma-landslide}), probably
        because a wall replaces the beach.
\end{itemize}

\paragraph{Plans.} In order of priority:
\begin{enumerate}[nosep]
  \item breaking: the transition just after the breaking point, the crest
        speed of the kinematic onset, and stabilised $k$--$\omega$
        \citep{larsen2018} for many-layer runs; the longshore current on the
        grid of \citet{ma2012};
  \item the landslide of Section~\ref{sec:ma-landslide} with its beach;
  \item volume-conserving wave-making zones and absorption of long waves;
  \item irregular wave generation and a boundary wavemaker;
  \item depth-integrated granular slides for landslide tsunamis;
  \item coupling to \code{interFoam} and \code{sedInterFoam} through relaxation
        zones, and hybrid meshes around structures.
\end{enumerate}
